\documentclass[11pt]{article}
\usepackage[margin=1in]{geometry}
\usepackage{amsmath,amssymb}
\usepackage[hidelinks]{hyperref}

\title{Literature-Implied Answer to the \(L^p\) Non-\(R\)-Schauder Basis Question}
\author{FA Banach run \texttt{fa\_banach\_001}, agent \texttt{agent\_lane\_12}}
\date{2026-06-20}

\begin{document}
\maketitle

\section*{Status}

This is a lightweight literature-implied answer packet. The source question is
answered affirmatively by combining a later general statement of
Arnold--Le Merdy with the classical fact that \(L^p([0,1])\), \(1<p<\infty\),
has an unconditional basis and is not Hilbert when \(p\neq2\).

\section*{Source Question}

Source paper:
Christoph Kriegler and Christian Le Merdy,
\emph{Tensor extension properties of \(C(K)\)-representations and applications
to unconditionality}, arXiv:0901.1025.

The question appears in Remark 5.6(2), page 22. After recalling that
Kalton--Lancien imply that \(L^p([0,1])\), \(1<p\neq2<\infty\), admits a
finite-dimensional Schauder decomposition which is not \(R\)-Schauder, the
authors write:
\begin{quote}
``However the question whether \(L^p([0,1])\) admits a Schauder basis which is
not \(R\)-Schauder, is apparently an open question.''
\end{quote}

\section*{Supporting Answer}

Supporting paper:
Loris Arnold and Christian Le Merdy,
\emph{New counterexamples on Ritt operators, sectorial operators and
\(R\)-boundedness}, arXiv:1812.11299.

On page 7, immediately after Corollary 0.4, Arnold--Le Merdy state:
\begin{quote}
``It follows from [4,6] that if \(X\) has an unconditional basis and \(X\) is
not isomorphic to a Hilbert space, then \(X\) has a Schauder basis which is not
\(R\)-Schauder.''
\end{quote}
Here [4] is Fackler's revisiting of the Kalton--Lancien theorem and [6] is
Kalton--Lancien, \emph{A solution to the problem of \(L^p\)-maximal regularity}.

\section*{Identification}

For \(1<p<\infty\), the Haar system is an unconditional Schauder basis of
\(L^p([0,1])\). For \(p\neq2\), \(L^p([0,1])\) is not isomorphic to a Hilbert
space. Applying the Arnold--Le Merdy statement with
\[
  X=L^p([0,1])
\]
therefore gives a Schauder basis of \(L^p([0,1])\) which is not
\(R\)-Schauder. This gives an affirmative answer to the full range asked in
Remark 5.6(2) of arXiv:0901.1025.

\section*{Scope and Provenance}

This is not a new proof from the run. The supporting authors do not appear to
explicitly identify Remark 5.6(2) of arXiv:0901.1025 as the question being
answered. The relation is an agent-identified implication from their general
statement.

The source's broader applications to \(C(K)\)-representations, \(H^\infty\)
calculus, and unconditionality are not reclassified here; only the specific
\(L^p([0,1])\) basis question in Remark 5.6(2) is cleared.

\begin{thebibliography}{9}

\bibitem{KrieglerLeMerdy2009}
C. Kriegler and C. Le Merdy,
\emph{Tensor extension properties of \(C(K)\)-representations and applications
to unconditionality},
arXiv:0901.1025.
\url{https://arxiv.org/abs/0901.1025}

\bibitem{ArnoldLeMerdy2018}
L. Arnold and C. Le Merdy,
\emph{New counterexamples on Ritt operators, sectorial operators and
\(R\)-boundedness},
arXiv:1812.11299.
\url{https://arxiv.org/abs/1812.11299}

\bibitem{Fackler2014}
S. Fackler,
\emph{The Kalton--Lancien theorem revisited: maximal regularity does not
extrapolate},
J. Funct. Anal. \textbf{266} (2014), 121--138.

\bibitem{KaltonLancien2000}
N. J. Kalton and G. Lancien,
\emph{A solution to the problem of \(L^p\)-maximal regularity},
Math. Z. \textbf{235} (2000), 559--568.

\end{thebibliography}

\end{document}
